\usepackage{amsmath,amssymb,amsthm,mathtools}
\usepackage{thmtools}
\usepackage[T1]{fontenc}
\usepackage{lmodern}
\usepackage{microtype}
\usepackage{enumitem}
\usepackage{booktabs}
\usepackage{array}
\usepackage{longtable}
\usepackage{etoolbox}
% Backport longtable v4.25's finite shrink at page breaks (LaTeX issue 1907).
% TeX Live 2025's mark extraction otherwise reports ignored primitive errors.
% Keep infinite stretch; change only the two shrink components in LT@output.
\makeatletter
\@ifpackagelater{longtable}{2025/12/22}{}{%
  \patchcmd{\LT@output}{\vss}{\vskip\z@\@plus1fil\@minus\normalbaselineskip}{}%
    {\PackageError{manuscript}{Cannot patch longtable page-break glue}{Update the longtable compatibility patch.}}%
  \patchcmd{\LT@output}{\vss}{\vskip\z@\@plus1fil\@minus\normalbaselineskip}{}%
    {\PackageError{manuscript}{Cannot patch longtable page-break glue}{Update the longtable compatibility patch.}}%
}
\makeatother
\usepackage[margin=1.05in]{geometry}
\PassOptionsToPackage{hyphens}{url}
\usepackage[hypertexnames=false,colorlinks=true,linkcolor=blue!50!black,citecolor=green!35!black,urlcolor=blue!60!black]{hyperref}
\usepackage[capitalise,noabbrev]{cleveref}
\renewcommand{\theHtable}{lt.\arabic{table}} % longtable anchors use \theHtable; keep them apart from float anchors
\Urlmuskip=0mu plus 1mu\relax % let long URLs stretch at their break points
\usepackage{tikz}
\usepackage{graphicx}
\renewcommand{\topfraction}{.95}
\renewcommand{\textfraction}{.05}
\renewcommand{\floatpagefraction}{.80}

\declaretheorem[name=Theorem,numberwithin=section,style=plain]{theorem}
\declaretheorem[name=Proposition,sibling=theorem,style=plain]{proposition}
\declaretheorem[name=Lemma,sibling=theorem,style=plain]{lemma}
\declaretheorem[name=Corollary,sibling=theorem,style=plain]{corollary}
\declaretheorem[name=Claim,sibling=theorem,style=plain]{claim}
\declaretheorem[name=Theorem,numbered=no,style=plain]{theorem*}
\declaretheorem[name=Definition,sibling=theorem,style=definition]{definition}
\declaretheorem[name=Construction,refname={Construction,Constructions},Refname={Construction,Constructions},sibling=theorem,style=definition]{construction}
\declaretheorem[name=Certificate,refname={Certificate,Certificates},Refname={Certificate,Certificates},sibling=theorem,style=definition]{certificate}
\declaretheorem[name=Hypothesis,refname={Hypothesis,Hypotheses},Refname={Hypothesis,Hypotheses},sibling=theorem,style=definition]{hypothesis}
\declaretheorem[name=Remark,sibling=theorem,style=remark]{remark}
\declaretheorem[name=Example,sibling=theorem,style=remark]{example}
\numberwithin{equation}{section}

\newcommand{\N}{\mathbb N}
\newcommand{\Z}{\mathbb Z}
\newcommand{\Q}{\mathbb Q}
\newcommand{\R}{\mathbb R}
\newcommand{\C}{\mathbb C}
\newcommand{\F}{\mathbb F}
\newcommand{\Fp}{\mathbb F_p}
\newcommand{\Zsix}{\Z/6\Z}

\newcommand{\E}{\mathbb E}
\renewcommand{\Pr}{\mathbb P}
\newcommand{\one}{\mathbf 1}
\newcommand{\ind}[1]{\mathbf 1_{\{#1\}}}
\newcommand{\abs}[1]{\left|#1\right|}
\newcommand{\norm}[1]{\left\|#1\right\|}
\newcommand{\ip}[2]{\left\langle #1,#2\right\rangle}
\newcommand{\opnorm}[1]{\left\|#1\right\|_{\mathrm{op}}}
\newcommand{\eps}{\varepsilon}
\newcommand{\dd}{\,\mathrm d}
\newcommand{\Var}{\operatorname{Var}}
\newcommand{\supp}{\operatorname{supp}}
\newcommand{\rad}{\operatorname{rad}}
\newcommand{\lcm}{\operatorname{lcm}}
\newcommand{\diag}{\operatorname{diag}}
\newcommand{\tr}{\operatorname{tr}}
\newcommand{\Unif}{\operatorname{Unif}}
\newcommand{\PSD}{\succeq}
\newcommand{\e}{e}            % e(t) = exp(2 pi i t) is written \ee
\newcommand{\ee}[1]{e\!\left(#1\right)}

\newcommand{\Qsq}[1][p]{Q_{#1}^{*}}           % nonzero squares mod p
\newcommand{\Ap}[1][p]{A_{#1}}                 % indicator of nonzero squares
\newcommand{\Om}[1][p]{\Omega_{#1}}            % strict-square pairs
\newcommand{\Xs}[1]{X_{#1}}                    % product of prime fields over a prime set

\newcommand{\Vslots}{V}                         % slot set (Z/6)^2 x {0,1}
\newcommand{\rot}{\mathsf c}                    % rotation on Z/6
\newcommand{\Rel}{\mathsf R}                    % relation y in {x, x+-1, c(x)}
\newcommand{\Pone}{P_1}
\newcommand{\Ptwo}{P_2}
\newcommand{\Pthree}{P_3}
\newcommand{\Pfour}{P_4}
\newcommand{\Mmod}{M}                           % small modulus 8 prod_{3<=p<P1} p
\newcommand{\Lift}{\mathcal L}                  % Fourier lift
\newcommand{\Yf}{\mathcal Y}                    % functional Y(N,A)
\newcommand{\Ym}{\mathfrak Y}                   % max over A
\newcommand{\aux}{\nu_{\mathrm{aux}}}           % auxiliary residue/interval law
\newcommand{\Mmap}{M_{\mathrm{map}}}
\newcommand{\Nmap}{N_{\mathrm{map}}}
\newcommand{\Nplus}{N_{+}}
\newcommand{\cut}[2]{(#1,#2)}                    % cut in coordinate #1 at position #2

\newcommand{\Pstar}{P^{*}}                      % start of the middle range (specified by the design)
\newcommand{\PpartII}{P_{\mathrm{II}}}          % start of Part II (equal to P* in the selected design)
\newcommand{\Pfive}{P_5}                        % start of the mixing tail = 2^65536 (sigma stage (c))
\newcommand{\Mzero}{M_0}                        % full physical modulus of the external index
\newcommand{\Kint}{K}                           % number of intervals of the external index
\newcommand{\Fc}{\mathfrak F}                   % coupled set of edge families, subset of Phi
\newcommand{\Dc}{\mathfrak D}                   % the transversal {(0,0),(1,1),(2,2)}
\newcommand{\ThetaF}[1][\Fc]{\Theta_{#1}}       % closing quantity Theta_F(a)
\newcommand{\Thetaone}{\Theta_{\mathrm{single}}} % centered selected-singleton closing quantity
\newcommand{\PF}[1][\Fc]{P_{#1}}                % union-bound polynomial of minimal F-obstructions
\newcommand{\MK}{M_{K}}                         % modulus moment M_K(2c)
\newcommand{\Qsqmod}[1]{\mathcal Q_{#1}}        % unit squares modulo an integer (Q_{M_0})
\newcommand{\Gom}{G_{\mathrm{om}}}              % one-mark field graph G_{s_row} on rows 0,1,2 (research: G_F)
\newcommand{\Gex}{G_{\mathrm{ex}}}              % exchangeable three-row graph (research: G_E)
\newcommand{\QQ}{\mathcal{Q}\!\mathcal{Q}}      % two-dimensional completion (p3_array Def. 4)
\newcommand{\QQphi}{\QQ^{\varphi}}              % type-weighted two-dimensional completion (p3_partII)
\newcommand{\Benv}{\mathsf B}                   % block constant B_0+B_1 (sum of the two positive-envelope constants)
\newcommand{\cplus}{c_{+}}                      % one-sided middle-row mismatch constant (= 2715)
\newcommand{\Rcut}{R_{\mathrm{cut}}}            % half-marginal density ratio of a cut
\newcommand{\Op}{\mathsf A_p}                   % normalised cut operator at the prime p
\newcommand{\gneg}[1]{\widehat g^{-}_{#1}}      % negative-side high-range constant, depth #1
\newcommand{\gpos}[1]{\widehat g^{+}_{#1}}      % positive-side high-range constant, depth #1
\newcommand{\Hmat}{\mathbf H}                   % star Gram (model) matrix
\newcommand{\Azero}{\mathbf A_{0}}              % E1 matrix chi(xi) eps_G Delta^E
\newcommand{\cperp}{\bar c_{\perp}}             % perp-block constant
\newcommand{\Lsc}{\mathsf L}                    % scale of the hybrid functional: N = K*Lsc (research: L)
\newcommand{\Hhi}{H_{\mathrm{hi}}}              % high-range exponent gamma_a W_3 + T_infty(a) in Theta (research: H)
\newcommand{\Jc}{\mathcal J}                    % random inclusion set of the coupling energy (research: J)
\newcommand{\cpk}[1][p]{\varpi_{#1}}            % coupling-kernel norm 1/(sqrt p - 1) (research: delta_p)

\newcommand{\Blk}{\mathcal B}                    % a block (atom) of primes; B alone is the bracket energy factor
\newcommand{\Nbr}{\mathsf N}                     % the reserved Gram bracket measure (never the scale N)
\newcommand{\Nbrz}{\mathsf N_0}                  % its record-difference-free part
\newcommand{\Jac}{\operatorname{Jac}}            % Jacobi sum
\newcommand{\Jall}{J_{\mathrm{all}}}             % the all-ones kernel
\newcommand{\Jproj}{\mathbb J}                   % the projection onto constants
\newcommand{\sq}{\operatorname{sq}}              % least square divisible by d
\newcommand{\mcl}{m_{\mathrm{cl}}}               % digit pairs of the square-clique prefix
\newcommand{\tside}{\tau_{\mathrm{side}}}        % direct-Gram side weight

\newcommand{\tlaw}{\mathsf t}                    % signed-law parameter 127/500 (never the repricing factor t)
\newcommand{\ulaw}{\mathsf u}                    % root-tilt parameter of a selected middle piece
\newcommand{\Hs}[1]{\mathsf H_{#1}}              % opposite-base-pair scores H_0, H_1, H_2
\newcommand{\Hbar}{\bar{\mathsf H}}              % their mean (the block density factor is \bar H_X)
\newcommand{\tagser}{\mathfrak h}                % tagged first-stage series h(d,v) and its majorant h_n^+
\newcommand{\Uinv}{U^{\mathrm{inv}}}             % assigned price of the invariant tags E_add+H
\newcommand{\Phat}{\hat P}                       % normalized Paley operator P/r
